What is a real number, a limit, a length? Mathematics used these words for two thousand years before anyone could
say what they meant, and the answers, found between 1858 and 1933, are among the most beautiful constructions in
the subject. This book tells that story and proves every step of it.

It builds the real line from nothing but sets, and shows that infinity comes in sizes. It frees limits from
coordinates with metric spaces and topology, and meets compactness and Baire's theorem, which says that most continuous functions have a derivative nowhere. A simple closed curve has an inside and an outside: that is the Jordan curve theorem. Then it
measures: Lebesgue's length, a set that has none, the integral that slices values instead of the axis, the volume of
the ball in every dimension, spaces of functions, densities as derivatives of measures, weak convergence, and
fractals whose dimension is not a whole number.

Every theorem is proved in full, every exercise has a hint and a reasoned solution. Written for the pleasure of seeing why things are true.
